\documentclass[12pt]{article}

% Idioma e codificação
\usepackage[brazil]{babel}
\usepackage[utf8]{inputenc}
\usepackage[T1]{fontenc}

% Matemática e listas
\usepackage{amsmath, amssymb}
\usepackage{enumitem}

% Layout e cores
\usepackage{geometry}
\usepackage{graphicx}
\usepackage{xcolor}

% Cores institucionais
\definecolor{maincolor}{RGB}{0,100,50}
\definecolor{forestgreen}{RGB}{0,100,50}

% Cabeçalho e rodapé
\usepackage{fancyhdr}

% Blocos
\usepackage[most]{tcolorbox}
\newtcolorbox{block}[1]{%
  colback=green!5,
  colframe=maincolor,
  coltitle=white,
  title=#1,
  fonttitle=\bfseries,
  boxrule=1.5pt,
  arc=3mm,
  left=6pt,
  right=6pt,
  top=6pt,
  bottom=6pt
}

% Margens
\geometry{margin=2.5cm}

% Fancyhdr
\pagestyle{fancy}
\fancyhf{}

\renewcommand{\headrulewidth}{2pt}
\renewcommand{\footrulewidth}{2pt}

\renewcommand{\headrule}{\hbox to\headwidth{%
  \color{forestgreen}\leaders\hrule height \headrulewidth\hfill}}

\renewcommand{\footrule}{\hbox to\headwidth{%
  \color{forestgreen}\leaders\hrule height \footrulewidth\hfill}}

% Altura do cabeçalho (corrigida)
\setlength{\headheight}{52pt}
\addtolength{\topmargin}{-20pt}

% Cabeçalho
\fancyhead[L]{%
  \includegraphics[height=1.2cm]{Universidade_de_Rio_Verde.jpg}
}

\fancyhead[R]{%
  \raggedleft
  \textcolor{forestgreen}{%
    \small
    \textbf{Lista de Exercícios 2}\\
    Me.~Sergio Souza Novak\\
    \textit{Engenharia da Qualidade e Confiabilidade}
    }
    }
    
    % Rodapé
    \fancyfoot[C]{\textcolor{forestgreen}{\thepage}}
    
    % =========================
    \begin{document}
    % =========================
    
    \vspace*{0.2cm}
% No preâmbulo, adicione este novo comando para caixas de resposta:
\newtcolorbox{resposta}[1]{%
  colback=white,
  colframe=forestgreen!30,
  arc=1mm,
  height=#1,
  width=\textwidth,
  boxrule=0.8pt
}

% Substitua o conteúdo da Parte II por este:


\begin{center}
  {\Large \textbf{Parte II: Métricas Quantitativas de Confiabilidade}}\\[0.2cm]
  \textit{Instruções: Apresente os cálculos e as unidades de medida para justificar suas respostas.}
\end{center}

\vspace{0.4cm}

\begin{block}{Fórmulas de Referência}
    \small
    \begin{itemize}
        \item $POFOD = \frac{\text{Nº de Falhas}}{\text{Nº de Solicitações}}$ \quad | \quad $ROCOF = \frac{\text{Nº de Falhas}}{\text{Tempo Total}}$
        \item $MTTF = \frac{1}{ROCOF}$ \quad | \quad $MTTR = \frac{\text{Tempo Total de Reparo}}{\text{Nº de Falhas}}$
        \item $AVAIL = \frac{MTTF}{MTTF + MTTR}$
    \end{itemize}
\end{block}

\vspace{0.4cm}

\textbf{1.} Um sistema de frenagem ferroviária realizou 2.500 testes e falhou em 5 deles. 
\begin{enumerate}[label=\alph*)]
    \item Calcule a \textbf{POFOD}.
    \item Se o requisito exige POFOD $\le 0,001$, o sistema está apto?
\end{enumerate}
\begin{resposta}{3.5cm}
\end{resposta}

\vspace{0.4cm}

\textbf{2.} Um servidor operou 720h com 6 interrupções. Determine a \textbf{ROCOF} e o \textbf{MTTF}.
\begin{resposta}{3.5cm}
\end{resposta}

\vspace{0.4cm}

\textbf{3.} Compare a Disponibilidade (\textbf{AVAIL}) dos sistemas abaixo:
\begin{itemize}
    \item \textbf{A:} Falha a cada 1.000h, reparo em 10h.
    \item \textbf{B:} Falha a cada 200h, reparo em 1h.
\end{itemize}
\begin{resposta}{4.5cm}
\end{resposta}

\vspace{0.4cm}

\textbf{4.} Para atingir 99,99\% de disponibilidade com $MTTF = 4.000h$, qual deve ser o \textbf{MTTR} máximo (em minutos)?
\begin{resposta}{4cm}
\end{resposta}


\textbf{5.} Um sistema de monitoramento ambiental opera 24/7. Em um trimestre (2.160 horas), o sistema apresentou um $MTTF$ de 500 horas. Calcule o número aproximado de falhas que ocorreram neste período.
\begin{resposta}{3.5cm}
\end{resposta}

\vspace{0.5cm}

\textbf{6.} Durante uma atualização de banco de dados, ocorreram 3 incidentes. O primeiro levou 15 min para resolver, o segundo 45 min e o terceiro 30 min. Calcule o \textbf{MTTR} médio desta operação.
\begin{resposta}{3.5cm}
\end{resposta}

\vspace{0.5cm}

\textbf{7.} Um e-commerce possui $AVAIL = 0,995$. Em um período de um ano (8.760 horas), quantas horas o sistema ficou, em média, indisponível para os clientes?
\begin{resposta}{3.5cm}
\end{resposta}

\vspace{0.5cm}

\textbf{8.} Explique a diferença prática entre utilizar \textbf{POFOD} e \textbf{ROCOF}. Em que tipo de sistema cada uma é mais adequada?
\begin{resposta}{4.5cm}
\end{resposta}

\vspace{0.8cm}
\vspace{0.5cm}

\textbf{9.} Na Engenharia de Software, os termos \textbf{Confiabilidade} e \textbf{Disponibilidade} são frequentemente confundidos, mas possuem focos distintos. Explique a diferença entre esses dois conceitos, abordando:
\begin{itemize}
    \item O que cada uma prioriza (continuidade vs. prontidão);
    \item Como um sistema pode ser altamente disponível, mesmo apresentando falhas frequentes.
\end{itemize}

\begin{resposta}{7cm}
\end{resposta}
\begin{center}
    \textit{“Boa atividade! Att. Me. Prof. Sergio Souza Novak.”}
\end{center}

% =========================
\end{document}
% =========================
